\subsection{Recommended proof and falsification sequence}

\begin{enumerate}[leftmargin=2em]
\item On endpoint-source paths, start from the proved
$\rho=\tau=q/5$ support scale and the moving-maximum upper bounds in
\cref{thm:path-moving-max-soft-upper}.  Seek an admission-specific energy,
normalized-error, or space--time potential proving
$K_{\rm face}=O(q^{-1})$, $T=O(q^{-2})$, and
$\mathfrak V_T=O(q^{-3})$; alternatively, close the audited short-path
candidate by deriving its literal entry coefficients, conditional-regime
margin, and anti-cancellation estimate.  Frontier-only logic is refuted, so
either route must retain the moving global correction.  Continue to track the
remaining-gain occupancy separately; the unweighted $q$-scaled shortcut is
already refuted.  Also avoid the two pointwise shortcuts refuted by
\cref{prop:path-monotone-correction-potential-fails}: the actual correction
need not decrease on a held face and need not be at most $q^{-1}$ times the
normalized face error with unit coefficient.  Also avoid treating
$\delta_t+q^{-1}e_t$ as a decreasing held-face bank: it passes stages
$6$--$7$ but fails on consecutive held stages $9$--$10$ by
\cref{prop:path-correction-error-bank-potential-fails}.  More generally, no
fixed $c\geq0$ makes $\delta_t+cq^{-1}e_t$ decrease on both pairs, by
\cref{prop:path-no-constant-coefficient-bank}.  These are pairwise local
constraints, not a bound on any global block horizon.  The first admissible
face-dependent test is the squared \textnormal{\textsc{TightPair--Schur}}
bank in \cref{prop:path-tight-pair-schur-bank}: debit the exact remaining
restricted-optimum drop only at the frozen-candidate reset, and never compare
the raw linear bank directly with objective shock.  Its immediate pointwise
extension is now refuted by
\cref{prop:path-tight-pair-local-chain-stop}: the reserve cannot pay either
adjacent fixed-face increase, even though it pays the reset.  The next upper
test in \cref{prop:path-tight-pair-recovery-block} shows that a literal
post-reset fixed-face telescope first repairs this one block at stage $13$,
but supplies no face-general or multi-admission rule.  The realized-credit
condition in \cref{prop:path-causal-two-admission-recovery} is the first
checkable face-general replacement: it uses $\delta^2$, debits every increase,
and credits only an already observed decrease or already computed
restricted-optimum drop.  It stays solvent across the consecutive stage-$4$
and stage-$8$ admissions because unused earlier credit supplies an exact
cross-state cancellation.  If that credit is discarded, stage $11$ remains a
ledger STOP and stage $12$ is the first local ledger GO.  The asymmetric
T-tree witness in \cref{prop:t-tree-causal-next-admission-stop} now refutes a
uniform recovery-before-the-next-admission rule: its stage-$5$ debt survives
the stage-$9$ admission and first closes at stage $14$.  The next upper task is
therefore narrower: prove an all-history solvency invariant or a structural
condition that supplies enough earlier credit, with every correction scan and
Schur query charged.  Alternatively, test a larger graph for debt that remains
negative through several subsequent admissions.  Do not promote either finite
trace to a global telescope.
\item Extend the proof to trees with bounded branching, isolating whether
simultaneous boundary admissions create an unavoidable $1/\alpha$ shock.
\item Test the zero-padded recurrence on adversarial spiders, stars of stars,
and lollipops.  Record peak volume, newly admitted volume, edge revisits,
$\nu(\ell_t)$, and objective shock separately.
\item If pointwise continuation fails, target the weaker charge
\cref{eq:new-volume-charge}; if that also fails, construct the corresponding
M-matrix counterexample and state the structural condition that restores it.
\end{enumerate}
